\documentclass[11pt]{article}
\usepackage[margin=1in]{geometry}
\usepackage[T1]{fontenc}
\usepackage{lmodern}
\usepackage{amsmath,amssymb,amsthm,mathtools}
\usepackage{booktabs,array}
\usepackage[numbers,sort&compress]{natbib}
\usepackage{microtype}
\usepackage[colorlinks=true,linkcolor=blue,citecolor=blue,urlcolor=blue]{hyperref}
\newcommand{\E}{\mathbb E}
\newcommand{\Var}{\operatorname{Var}}
\newcommand{\Cov}{\operatorname{Cov}}
\newcommand{\N}{\mathbb Z_{\geq0}}
\newcommand{\ind}{\mathbf 1}
\newcommand{\dd}{\,\mathrm dt}
\newcommand{\norm}[1]{\left\lVert #1\right\rVert}
\DeclareMathOperator{\supp}{supp}
\newtheorem{theorem}{Theorem}
\newtheorem{proposition}{Proposition}
\newtheorem{lemma}{Lemma}
\title{A local noise tradeoff for two cross-regulated genes\\
with unequal finite RNA lifetimes}
\author{Xiaojian Zeng\\
\small PBC School of Finance, Tsinghua University, 100190\\
\small \texttt{xiaojianzeng@foxmail.com}}
\date{}
\hypersetup{pdftitle={A local noise tradeoff for two cross-regulated genes with unequal finite RNA lifetimes},pdfauthor={Xiaojian Zeng}}
\begin{document}
\maketitle
\begin{abstract}
Eliminating short-lived RNA gives an instantaneous protein-burst model,
but a noise constraint in that limit need not describe finite RNA lifetimes.
We establish a computer-assisted local constraint for two genes with explicit
RNA, catalytic translation and protein degradation. In protein-lifetime units,
the constitutive transcription and RNA-decay rates are respectively $(1,1)$
and $(2,2)$, while both translation and protein-decay coefficients are one.
Each parent has RNA and protein mean one. Cross-transcription responses may
be any bounded real functions of the other protein count, including
nonmonotone functions and responses confined to arbitrarily rare counts.
An infinite-state return-operator bound, obtained from weighted generator
estimates and exact finite Riesz residual certificates, is combined with a
controlled cubic covariance remainder. For each fixed nonconstant response
pair, an explicit positive feedback interval forces at least one stationary
protein Fano factor above its own mean-matched finite-lifetime baseline,
$3/2$ or $4/3$. Constant responses leave both baselines unchanged.
The interval depends on the response shapes. The result does not settle
strong feedback, other kinetic parameters or networks with more genes.
\end{abstract}

\section{Finite RNA memory and output-noise constraints}

At low molecule counts, a gene-regulatory circuit is naturally represented
by stochastic reaction events. Noise control must then be compared with
an uncontrolled system at the same stationary mean, rather than with an
arbitrarily chosen production rate. For explicit RNA and catalytic
translation, the constitutive protein Fano factor is
\begin{equation}
 B_i=1+\frac{\beta_i}{\delta_i+\gamma_i},
 \label{eq:baseline}
\end{equation}
where $\delta_i$ and $\gamma_i$ are RNA and protein decay coefficients
and $\beta_i$ is the translation coefficient
\citep{Thattai2001,Shahrezaei2008}. This finite-lifetime baseline differs
from the instantaneous-burst value $1+\beta_i/\delta_i$.

Yan \emph{et al.} formulated noise-tradeoff conjectures for reaction
networks and established analytical constraints in linear-noise and
separated-time regimes \citep{Yan2019}. Stationary information-flow
arguments subsequently supplied exact constraints for unit-birth
components relative to their Poisson baselines
\citep{Ripsman2025,AndersonUnit2026}. An exact tradeoff also holds for
instantaneous positive geometric protein bursts, while unrestricted
nongeometric burst laws admit counterexamples \citep{AndersonBurst2026}.
The latter work identifies finite-lived production intermediates as a
remaining direction. A protein-only burst generator omits the RNA memory
retained here. Likewise, a constraint on all RNA and protein coordinates
relative to one does not by itself compare only the proteins with
Eq.~\eqref{eq:baseline}.

We consider a restricted, exactly specified part of that direction:
two genes with unequal finite RNA lifetimes and low constitutive mean
counts. The contribution is a local output-noise obstruction valid for
every bounded response shape on the complete count lattice. The response
functions are not linearized at their means. The proof combines exact
stationary identities, a computer-assisted bound on the complete
response-function space, and an explicit finite-amplitude remainder.
The fixed kinetics and response-dependent amplitude are essential parts
of the statement.

\section{Reaction model and the exact comparison}
\label{sec:model}

Let $(M_1,P_1,M_2,P_2)\in\N^4$ denote RNA and protein counts. For
$i=1,2$, with $j\ne i$, the only channels are
\begin{align}
 M_i&\longrightarrow M_i+1 &&\text{at rate } f_i(P_j),&
 M_i&\longrightarrow M_i-1 &&\text{at rate } \delta_i M_i,
 \nonumber\\
 P_i&\longrightarrow P_i+1 &&\text{at rate } \beta_i M_i,&
 P_i&\longrightarrow P_i-1 &&\text{at rate } \gamma_i P_i.
 \label{eq:reactions}
\end{align}
Translation leaves RNA unchanged. Death rates vanish at zero, and there
is no upper count boundary. Transcription depends only on the other
protein: there is no direct protein self-regulation, promoter state,
degradation feedback, extrinsic environment or additional controller.
No detailed-balance assumption is imposed.

Choose a common time unit $t_0$ equal to the protein mean lifetime.
All times and rates below are their numerical values in units $t_0$
and $t_0^{-1}$. Counts, response functions and $\epsilon$ are dimensionless;
$a,b$ are numerical transcription gains in the same rate unit.
The fixed parameters are
\begin{equation}
 (\delta_1,\beta_1,\gamma_1)=(1,1,1),\qquad
 (\delta_2,\beta_2,\gamma_2)=(2,1,1).
 \label{eq:kinetics}
\end{equation}
For arbitrary functions $h_i:\N\to[-1,1]$ and gains $a,b>0$, set
\begin{equation}
 f_1(p)=1+\epsilon a h_2(p),\qquad
 f_2(p)=2-\epsilon b h_1(p),\qquad g=\max(a,b).
 \label{eq:feedback}
\end{equation}
The sign convention places no monotonicity constraint on $h_i$.
For $\epsilon g<1$, both rates are positive and bounded at every count.

These conditions define an admissible infinite-state process. Its total
jump rate grows at most linearly, so it is nonexplosive. Positive
immigration, translation and death allow communication between all
states. For bounded transcription the Foster function
\begin{equation}
 \Phi=1+\sum_{i=1}^2\left[\left(1+\frac{1}{\delta_i}\right)M_i+P_i\right]
\end{equation}
has negative linear drift outside a finite set. Polynomial powers of
$\Phi$ have the corresponding negative leading drift, because jump
increments are bounded and jump rates grow at most linearly. This gives
positive recurrence, a unique stationary law $\pi_\epsilon$ and the
polynomial moments needed below. Protein means are strictly positive.

Write $\mu_i=\E_\epsilon P_i$ and
$F_i=\Var_\epsilon(P_i)/\mu_i$. The comparison gene has exactly the same
$\delta_i,\beta_i,\gamma_i$ and constant transcription chosen to match
$\mu_i$. Its Fano factor is Eq.~\eqref{eq:baseline}, independent of that
constant rate. Thus our baselines are
\begin{equation}
 B_1=\frac32,\qquad B_2=\frac43.
 \label{eq:fixedbaselines}
\end{equation}
They refer to the actual regulated means, not only to the parent means.

\begin{proposition}[Exact finite-lifetime noise identity]
\label{prop:identity}
For the bounded reaction model~\eqref{eq:reactions}, with any finite
positive kinetic coefficients,
\begin{equation}
 F_i-B_i=
 \frac{\beta_i\,\Cov_\epsilon(Z_i,f_i)}
 {\gamma_i\mu_i(\delta_i+\gamma_i)},\qquad
 Z_i=\frac{\beta_i}{\delta_i}M_i+P_i.
 \label{eq:fanoidentity}
\end{equation}
\end{proposition}
\begin{proof}
Stationary generator tests on $M_i,P_i,M_i^2,M_iP_i,P_i^2$ give
\begin{align}
 \delta_i\E M_i&=\E f_i,&
 \gamma_i\E P_i&=\beta_i\E M_i,\nonumber\\
 \Var(M_i)&=\E M_i+\frac{\Cov(M_i,f_i)}{\delta_i},\nonumber\\
 \Cov(M_i,P_i)&=
 \frac{\beta_i\Var(M_i)+\Cov(P_i,f_i)}{\delta_i+\gamma_i},\nonumber\\
 \Var(P_i)&=\E P_i+\frac{\beta_i}{\gamma_i}\Cov(M_i,P_i).
 \label{eq:momentidentity}
\end{align}
Substitution proves Eq.~\eqref{eq:fanoidentity}. The polynomial moment
bounds justify the generator tests by localization.
\end{proof}
In the fixed units~\eqref{eq:kinetics}, $Z_1=M_1+P_1$,
$Z_2=M_2/2+P_2$, and
\begin{equation}
 F_i-B_i=\frac{\Cov_\epsilon(Z_i,f_i)}{(\delta_i+1)\mu_i}.
 \label{eq:fixedidentity}
\end{equation}
In particular, the sign of the excess Fano factor is the sign of the
corresponding covariance.

\section{The local result}
\label{sec:result}

At $\epsilon=0$, the two genes are independent constitutive parents,
with $(c,\delta,\beta,\gamma)=(1,1,1,1)$ and $(2,2,1,1)$.
Both RNA and protein means are one. Let $\mathcal T_i(t)$ be the
single-parent backward semigroup and $\E_0$ expectation under the
independent stationary parents. Define
\begin{align}
 A_i&=\Cov_0(M_i,h_i(P_i)),&
 H_i&=\Cov_0(P_i,h_i(P_i)),\nonumber\\
 C_i(t)&=\Cov_0(h_i(P_i(0)),h_i(P_i(t))),\nonumber\\
 D_i(t)&=\E_0\big[(\mathcal T_i(t)h_i)(M_i+1,P_i)
                      -(\mathcal T_i(t)h_i)(M_i,P_i)\big].
 \label{eq:parentforms}
\end{align}
The last expression is an extra-RNA response, not a simultaneous
RNA/protein jump. Put
\begin{equation}
 r_1(t)=(1+t)e^{-t},\qquad
 r_2(t)=e^{-t}-\tfrac12e^{-2t},
\end{equation}
and define the exact injection and return forms
\begin{align}
 U_1&=\int_0^\infty r_1(t)C_2(t)\dd,&
 U_2&=\int_0^\infty r_2(t)C_1(t)\dd,\nonumber\\
 V_1&=A_1\int_0^\infty r_1(t)D_2(t)\dd+
             H_1\int_0^\infty e^{-t}D_2(t)\dd,\nonumber\\
 V_2&=A_2\int_0^\infty r_2(t)D_1(t)\dd+
             H_2\int_0^\infty e^{-t}D_1(t)\dd.
 \label{eq:UV}
\end{align}
Global additive constants in the responses do not affect these forms.
The integrals are finite for bounded $h_i$.

\begin{theorem}[Local finite-RNA output-noise tradeoff]
\label{thm:main}
For the fixed parents above, every pair of bounded real whole-lattice
responses satisfies
\begin{equation}
 |V_1|,\ |V_2|\ \leq\ \frac45\sqrt{U_1U_2}.
 \label{eq:contraction}
\end{equation}
For the physical family~\eqref{eq:feedback}, define
$S=a^2U_1+b^2U_2$. If at least one $h_i$ is nonconstant, then $S>0$.
For every
\begin{equation}
 0<\epsilon<
 \min\left\{\frac1g,\frac{S}{4480g^3}\right\},
 \label{eq:epsilon}
\end{equation}
the admissible stationary network has
\begin{equation}
 F_1>\frac32\quad\text{or}\quad F_2>\frac43.
 \label{eq:tradeoff}
\end{equation}
If both response functions are constant, $F_1=3/2$ and $F_2=4/3$.
\end{theorem}

The operator inequality~\eqref{eq:contraction} is computer-assisted:
its finite matrices, rational witnesses and infinite-state error bounds
are described in Sec.~\ref{sec:contraction} and Appendix~\ref{app:certificate}.
The statement has no finite-plateau, smoothness or monotonicity restriction.
It has a response-dependent amplitude: nearly constant or rare-tail
responses can make $S$ arbitrarily small.

\subsection{Finite-amplitude response and deduction}

The independent-product Poisson equation for
$(Z_i-\E_0 Z_i)(h_j-\E_0 h_j)$ gives the quadratic covariance terms
in Eq.~\eqref{eq:UV}. A stationary coupling and weighted increment
estimate supply the following finite remainder, proved in
Appendix~\ref{app:remainder}:
\begin{align}
 \left|\Cov_\epsilon(Z_1,f_1)
       -\epsilon^2(a^2U_1-abV_1)\right|
 &\leq448\epsilon^3g^3,\nonumber\\
 \left|\Cov_\epsilon(Z_2,f_2)
       -\epsilon^2(b^2U_2-abV_2)\right|
 &\leq448\epsilon^3g^3 .
 \label{eq:remainder}
\end{align}
These are whole-lattice inequalities for $\epsilon g<1$; stationary
differentiability is not assumed.

Using Eq.~\eqref{eq:contraction} and
$2ab\sqrt{U_1U_2}\leq S$,
\begin{equation}
 S-ab(V_1+V_2)\geq\frac15S.
\end{equation}
Hence
\begin{equation}
 \Cov_\epsilon(Z_1,f_1)+\Cov_\epsilon(Z_2,f_2)
 \geq \epsilon^2\frac S5-896\epsilon^3g^3>0
\end{equation}
under Eq.~\eqref{eq:epsilon}. At least one covariance is positive, and
Eq.~\eqref{eq:fixedidentity} proves Eq.~\eqref{eq:tradeoff}.
Strict positivity of $S$ follows from the weighted estimates below.
When both responses are constant, the genes remain independent
constitutive systems with changed means, so their Fano factors remain
Eq.~\eqref{eq:fixedbaselines}.

\section{Controlling all response shapes}
\label{sec:contraction}

\subsection{Injection variance and a weighted generator estimate}

For the parent with RNA decay $d\in\{1,2\}$, let $\pi_d$ denote its
stationary law and $L_d$ its full backward generator. Its protein
marginal has full support. For a bounded $h$, write
$\widetilde h=h-\E_{\pi_d}h$ and
\begin{equation}
 W_{4,d}(h)=\E_{\pi_d}
       \frac{\widetilde h(P)^2}{(1+P)^4}.
\end{equation}
The scalar stationary covariance is even in time, although the parent
process is not reversible. The impulse-filter covariance identity gives
\begin{align}
 U_1&=2\Var_{\pi_2}
       \left(\int_0^\infty te^{-t}\widetilde h_2(P_t)\dd\right),\nonumber\\
 U_2&=3\Var_{\pi_1}
       \left(\int_0^\infty(e^{-t}-e^{-2t})
                         \widetilde h_1(P_t)\dd\right).
 \label{eq:filtervariance}
\end{align}
One can verify these identities by integrating the autocorrelations of
the displayed kernels. Bounded responses justify Fubini. Conditioning on
$(M_0,P_0)$ and applying Jensen yields
\begin{align}
 U_1&\geq2\norm{(1-L_2)^{-2}\widetilde h_2}_2^2,\nonumber\\
 U_2&\geq3\norm{(1-L_1)^{-1}(2-L_1)^{-1}\widetilde h_1}_2^2.
 \label{eq:reslower}
\end{align}
All norms are under the appropriate full invariant law.

Let $q_d(m,p)=d+(d+1)m+p$ be the total outgoing rate and
$w_d=1+q_d$. Neighboring states have $w_d$ ratios between $1/2$ and $2$.
Stationary incoming flux is $\pi_d(y)q_d(y)$. Outgoing-jump
Cauchy--Schwarz, followed by this incoming-flux identity, gives
\begin{equation}
 \norm{w_d^{-j}L_du}_2^2
 \leq2(1+2^{2j-1})\norm{w_d^{-(j-1)}u}_2^2,
 \qquad j=1,2.
 \label{eq:genbound}
\end{equation}
For the incoming term, the weight ratio costs at most $2^{2j-1}$;
the outgoing diagonal term costs at most one. The factor two comes
from $(u(y)-u(x))^2\leq2u(y)^2+2u(x)^2$.
This uses nonnegative countable sums, not detailed balance.
Applying Eq.~\eqref{eq:genbound} twice gives
\begin{align}
 \norm{w_d^{-2}(1-L_d)^2v}_2
 &\leq(1+2\sqrt6+\sqrt{108})\norm{v}_2<17\norm{v}_2,\nonumber\\
 \norm{w_d^{-2}(1-L_d)(2-L_d)v}_2
 &\leq(2+3\sqrt6+\sqrt{108})\norm{v}_2<20\norm{v}_2.
 \label{eq:productbound}
\end{align}
The resolvent vectors in Eq.~\eqref{eq:reslower} have
bounded first and second generator powers, so their backward resolvent
equations justify these operations without a finite-support-core
assumption.

The stationary cluster construction in Appendix~\ref{app:parent}
implies $\E[M\mid P=p]\leq p+1$. Thus
$\E[w_d\mid P=p]\leq2(d+1)(p+1)$. Conditional Jensen for $x^{-4}$,
Eq.~\eqref{eq:reslower} and Eq.~\eqref{eq:productbound} give
\begin{equation}
 U_1\geq\frac{W_{4,2}(h_2)}{187272},
 \qquad
 U_2\geq\frac{3W_{4,1}(h_1)}{102400}.
 \label{eq:coercivity}
\end{equation}
The denominators follow from
$2/(17^2\,6^4)=1/187272$ and $3/(20^2\,4^4)=3/102400$.
Every nonconstant bounded response therefore has strictly positive
injection. The weight is deliberately weak; no stronger unproved
coercivity estimate is needed.

\subsection{Finite certificates and omitted functionals}

Let $\mathcal H_d$ be the completion of bounded protein functions
modulo constants under injection norm $U_2$ for $d=1$, or $U_1$
for $d=2$. Write the lower constants in Eq.~\eqref{eq:coercivity} as
$\alpha_1=3/102400$, $\alpha_2=1/187272$.
For a zero-sum coefficient functional $\ell(h)=\sum_p\ell_ph(p)$,
\begin{equation}
 \norm{\ell}_{\mathcal H_d^*}^2
 \leq\frac1{\alpha_d}
       \sum_{p\geq0}\frac{(p+1)^4\ell_p^2}{\pi_d(P=p)}.
 \label{eq:dual}
\end{equation}
This follows by centering $h$ and applying Cauchy--Schwarz.

The return forms have rank two. With
$D_{g,i}=\int e^{-t}D_i(t)\dd$ and
$D_{r,i}=\int r_i(t)D_i(t)\dd$,
\begin{equation}
 V_1=A_1D_{r,2}+H_1D_{g,2},\qquad
 V_2=D_{r,1}A_2+D_{g,1}H_2.
 \label{eq:ranktwo}
\end{equation}
Appendix~\ref{app:parent} provides whole-lattice coefficient bounds
$|\ell_p|\leq C(p+1)\pi_d(P=p)$ for each of these functionals.
They are consequently continuous in $\mathcal H_d$.

Let $\mathcal S_n$ consist of response functions constant for $p\geq n$,
modulo global constants, and let $\mathcal P_{d,n}$ be its
injection-orthogonal projection. A full threshold category is included;
this is not a cap on the stochastic process. Exact interval Gram matrices
and integer congruence certificates at $n=64$ establish
\begin{equation}
 |V_i(\mathcal P_{1,n}h_1,\mathcal P_{2,n}h_2)|
 \leq\frac34\sqrt{U_2(h_1)U_1(h_2)},\qquad n\leq64.
 \label{eq:finitebound}
\end{equation}
The parent resolvents are approximated on $0\leq M\leq48$,
$0\leq P\leq128$ only to construct witnesses. Their errors under the
true invariant law, including the infinite tail, enter every matrix
enclosure. Appendix~\ref{app:certificate} specifies that error control.

For each return functional, a rational vector $v\in\mathcal S_n$
approximates its Riesz representative. Write its residual as
$e=\ell-U_d(v,\cdot)$, where $U_d$ denotes the corresponding polarized
injection form. In the scaled basis
$s_k\ind_{\{P=k\}}$ for $1\leq k<n$ and
$s_n\ind_{\{P\geq n\}}$, let $\rho_k$ enclose $|e(b_k)|$.
The count-zero coefficient is bounded by $\sum_k\rho_k/s_k$,
because $e$ annihilates constants and the last basis function covers
the entire tail. The head dual bound is
\begin{equation}
 H=\frac{(\sum_{k=1}^n\rho_k/s_k)^2}{\pi_{d,0}^{-}}
    +\sum_{p=1}^{n-1}
       \frac{(p+1)^4(\rho_p/s_p)^2}{\pi_{d,p}^{-}},
 \label{eq:head}
\end{equation}
where $\pi_{d,p}^{-}>0$ is a certified lower probability bound.

If $V\geq\sup_p|v(p)-\E v|$, the coefficient of the injection
covariance operator obeys
$|(U_dv)_p|\leq c_dV\pi_{d,p}$, with $c_1=3/4$ and $c_2=2$.
The omitted dual contribution is therefore at most
\begin{align}
 T&=C^2T_6+2Cc_dVT_5+(c_dV)^2T_4,\nonumber\\
 T_k&\geq\E_{\pi_d}[(P+1)^k\ind_{\{P\geq n\}}],\qquad k=4,5,6.
 \label{eq:tail}
\end{align}
Exact full moments minus lower probability prefixes enclose these
positive tails. Thus
\begin{equation}
 \Delta_\ell=\sqrt{(H+T)/\alpha_d}
 \label{eq:rieszerror}
\end{equation}
is an upper bound on the full dual residual and on the norm of the
omitted functional $\ell-\ell\mathcal P_{d,n}$.
An upper bound for the projected functional is
$m_\ell=\norm v_{U_d}+\Delta_\ell$.
For a factor pair $\ell,k$ in Eq.~\eqref{eq:ranktwo}, its omitted
operator norm is bounded by
\begin{equation}
 \Delta_\ell m_k+m_\ell\Delta_k+\Delta_\ell\Delta_k.
 \label{eq:rankerror}
\end{equation}
Adding the two rank-one errors to $3/4$ bounds the full return norm.
No finite projection is assumed to exhaust the response space.

\input{manuscript/certificate-table.tex}

The decisive $n=48$ exact bounds in Table~\ref{tab:certificate} are
strictly below $4/5$. They prove Eq.~\eqref{eq:contraction} for all
bounded response shapes. The smaller projection gives valid bounds
below one but insufficient to certify $4/5$; it is included to show
the effect of the residual enclosure. Floating solves propose witnesses;
exact rational residuals, probability intervals, moment tails and integer
congruence margins decide all displayed inequalities.

\section{Physical interpretation and limitations}

The positive injection forms measure fluctuations introduced by the
other gene's transcription signal after RNA-to-protein filtering.
The return terms measure how the loop acts back on that noise.
At the stated low-copy parents, the full return response is too weak
to offset both injection costs locally, even if the response is
nonlinear, nonmonotone or concentrated on very rare protein counts.
Finite RNA memory is retained throughout.

The result restricts a design route rather than asserting a universal
noise law. It does not resolve simultaneous protein-noise suppression
for arbitrary kinetic parameters, strong feedback or more than two genes.
It also supplies no response-independent positive $\epsilon$ radius.
The conservative constants are sufficient bounds, not optimal feedback
thresholds or experimental predictions. Other channels, including
promoter switching, resource coupling, sequestration and degradation
control, lie outside Eq.~\eqref{eq:reactions}. A comparison with a burst
baseline at finite RNA lifetime would answer a different question.

The distinction from prior exact unit-birth and geometric-burst results
is the explicit finite-RNA, output-only, mean-matched comparison.
The literature comparison here is scoped to those primary models and
the analytical regimes of Ref.~\citep{Yan2019}; it is not a claim to
have excluded every representation of a related result.

\section*{Data and computational methods}
The companion source package contains the exact matrices, probability
intervals, rational witnesses, residual bounds and certificate-replay code.
The enclosure construction used Python with NumPy and SciPy for
proposing sparse-resolvent and Riesz approximants; all final certificate
signs use integer or \texttt{Fraction} arithmetic. A standard-library
checker replays the stored contraction and residual certificates.
Appendices below specify the underlying infinite-state arguments.
There are no experimental data. The local package has no public
repository identifier at this draft stage.

OpenAI Codex was used substantively in exploring the argument,
literature synthesis, derivations, code generation, automated
computational checks and manuscript drafting. The automated checks
describe reproducible calculations, not completed human verification.
Human scientific review and approval remain pending for this local
draft. The accompanying author handoff records the disclosure details
and model-version provenance that require confirmation before submission.

\appendix
\input{manuscript/appendices.tex}
\bibliographystyle{unsrtnat}
\bibliography{manuscript/references}
\end{document}
